\subsection{RQ1.1: Coarse Boundary Selection}
\label{sec:analysis-rq11}

RQ1.1 asks which coarse ResNet-18 stage boundaries provide suitable
two-part microservice decompositions under controlled local execution.
The analytical contribution of this section is not the per-condition
latency itself, which is reported in \cref{tab:rq11-results}, but the
explanation of why the four coarse boundaries are not interchangeable
and which of them is the strongest carry-forward candidate. The
RQ1.1 results support a single coherent reading: the cost of one
service boundary is a joint function of the activation transferred at
that boundary and of how the remaining computation is distributed
between the two services. This matches the framing used in the prior
split-inference literature, which treats partition placement as a
computation--communication trade-off rather than a depth-only
decision~\cite{neurosurgeon,dnnsplit}. The four stage boundaries of
ResNet-18~\cite{resnet} provide a structurally clean coarse sweep, but
the RQ1.1 data make clear that this structural cleanliness is a
necessary, not sufficient, condition for a suitable decomposition.

\paragraph{All four splits are slower than monolithic.}
Under the controlled local conditions used here, no two-part
decomposition improves raw end-to-end latency relative to monolithic
inference: the four splits add between \(+2.4\%\) and \(+4.7\%\) mean
overhead in \cref{tab:rq11-results}. The RQ1.1 result is therefore
diagnostic rather than optimisation-oriented; it identifies the
least-bad boundary and explains the gradient of costs across the
sweep, but it does not support a claim that local two-service
decomposition is a raw-latency improvement. This is consistent with
the broader split-inference literature, where useful latency wins
typically depend on heterogeneous compute or bandwidth conditions that
are deliberately absent from RQ1.1~\cite{neurosurgeon,dnnsplit}.

\paragraph{The earliest boundary is penalised by activation expansion.}
The split after \texttt{layer1} transferred the largest activation
tensor (\(802{,}816\)~B, \(784\)~KiB) and had the largest mean
boundary-crossing interval at \(53.81\)~ms (\cref{tab:rq11-results}).
A pointed comparison is informative: the raw input tensor is
\(1\times3\times224\times224\) FP32, i.e.\ \(602{,}112\)~B, so the
intermediate carried at this boundary is approximately \(1.33\times\)
the raw input. The earliest coarse stage therefore \emph{expands} the
byte budget that the service interface must carry, which is the
cleanest in-thesis explanation for why the earliest split is the most
expensive. Prior work on split inference for CNN backbones reports the
same qualitative pattern --- without representation compression, the
early-stage feature maps can exceed the raw input
size~\cite{bottlenet_pp,deeperthings}. The boundary-crossing interval
reported here is, however, a composite measurement that includes
remote service-B compute and serialisation/deserialisation in addition
to transport (\cref{sec:rq11-results}), so it should not be read as a
pure network-transfer cost.

\paragraph{The latest boundary is penalised by compute degeneracy.}
The split after \texttt{layer4} shows the symmetric failure mode. It
transferred the smallest activation tensor (\(98\)~KiB) and had the
lowest boundary-crossing interval (\(2.22\)~ms), but it left only
\(0.486\)~ms of mean compute on the downstream service, approximately
\(0.62\%\) of total split compute. This is well below the predeclared
\(10\%\) compute-degeneracy threshold (\cref{sec:methods-carryforward})
and the condition is correspondingly flagged in
\texttt{carry\_forward.json}. The implication is methodologically
important: minimising the transferred activation alone can produce a
\emph{structurally weak} microservice decomposition that reduces
transfer burden by collapsing the second service into little more
than the classification head. The same pattern is described in the
distributed-CNN literature, which observes that very late splits
imbalance compute across the partition~\cite{deeperthings,neurosurgeon}.
RQ1.1 therefore retains \texttt{split\_after\_layer4} as evidence for
the transfer-size gradient, not as a credible carry-forward candidate.

\paragraph{The carry-forward region is mid-to-late.}
With the earliest and latest boundaries excluded for the reasons
above, the strongest RQ1.1 candidates are the two mid-to-late
boundaries. The raw-fastest non-degenerate split is
\texttt{split\_after\_layer3} at \(79.08\)~ms, which is selected as
the main carry-forward candidate by the frozen
\texttt{carry\_forward.json}; \texttt{split\_after\_layer2} is
retained as the reference at \(80.12\)~ms. The difference between
these two is modest (approximately \(1.04\)~ms, around \(1.3\%\) of
the slower mean) and should not be read as a general ranking of
\texttt{layer3} over \texttt{layer2} across all environments. The
defensible conclusion is narrower: under RQ1.1's controlled local
conditions, the credible carry-forward region lies around
\texttt{layer2}--\texttt{layer3}, where activation transfer has
already fallen substantially but downstream compute remains
meaningful. This region is what RQ1.2 then refines.

\paragraph{Validation.}
Three pieces of evidence support treating the within-run ordering and
the carry-forward decision as credible for the measured environment.
First, functional parity holds: both local and gRPC parity checks
report \(\texttt{max\_abs\_diff} = 0.0\) at the predeclared tolerance
of \(1.0\times10^{-5}\) across all four split conditions
(\texttt{parity\_validation.json}), so latency differences are not
confounded with numerical drift. Second, round-level behaviour is
stable: the per-condition round coefficient of variation lies between
\(0.002\) (monolithic) and \(0.011\) (\texttt{split\_after\_layer2})
in \texttt{round\_consistency.json}, in line with the
benchmarking-practice guidance to report round-level consistency
rather than mean alone~\cite{benchmarking_sc15}. Third, host controls
(single physical core, thread counts pinned to one, performance
governor, turbo disabled) are applied symmetrically across all
conditions (\cref{tab:rq11-config}), so the comparison is internally
fair. These three legs of the validation argument support the
within-run reading; they do not establish host-independent generality,
because RQ1.1 uses one host, one build, one fixed input seed, CPU-only
batch-one inference, and localhost gRPC, and cross-host or cross-seed
replication is not in scope for this stage.

\paragraph{Evaluation.}
The evaluative reading of RQ1.1 has two facets. Against the
monolithic baseline, every split is measurably slower, with relative
overheads of \(+4.7\%\), \(+3.7\%\), \(+2.4\%\) and \(+3.0\%\) for
splits after \texttt{layer1}, \texttt{layer2}, \texttt{layer3} and
\texttt{layer4} respectively (\cref{tab:rq11-results}). RQ1.1 is
therefore not a result about whether decomposition is faster; it is a
result about which decomposition is least costly while preserving a
meaningful microservice. Within the split set, the carry-forward rule
produces a defensible main--reference pair whose within-run per-iteration
\(95\%\) confidence intervals do not overlap
(\cref{tab:rq11-results}). That separation is a within-run property
and is not a host-independent uncertainty estimate, so it justifies
carrying both \texttt{split\_after\_layer3} and
\texttt{split\_after\_layer2} forward to RQ1.2 rather than committing
to a single boundary at this stage.

\paragraph{Answer to RQ1.1.}
RQ1.1 provides direct but scoped evidence that the most suitable
coarse two-part ResNet-18 microservice decompositions under controlled
local execution are not the earliest or latest stage boundaries. The
earliest boundary is penalised because the layer-1 activation expands
the byte budget the boundary must carry, exceeding the raw input size
in this run. The latest boundary is penalised because it satisfies
the activation-size criterion only by reducing the downstream service
to roughly the classification head, falling below the predeclared
\(10\%\) compute-degeneracy threshold. The strongest carry-forward
candidate from this frozen run is \texttt{split\_after\_layer3}, with
\texttt{split\_after\_layer2} retained as the reference candidate.
Suitability under RQ1.1 conditions is therefore determined by the
\emph{joint} balance of latency, boundary-crossing cost, activation
size and compute distribution, not by any single dimension --- the
two-dimensional reading developed further in
\cref{sec:analysis-rq13}.
